%---------- Inleiding ---------------------------------------------------------


\section{Inleiding}%
\label{sec:inleiding}

Professionele voetbalclubs streven voortdurend naar sportieve en competitieve vooruitgang. 
In dit streven wordt steeds vaker gebruikgemaakt van data en technologie om spelersprestaties te monitoren, te evalueren en te optimaliseren. 
Door het gebruik van GPS-systemen, wearables en geavanceerde wedstrijdstatistieken beschikken clubs vandaag over een grote hoeveelheid objectieve fysieke prestatiegegevens.

Deze bachelorproef focust zich op het analyseren van fysieke spelersattributen binnen een wedstrijdcontext en onderzoekt hoe machine learning kan bijdragen aan het identificeren van fysiek relevante spelersprofielen. 
Door fysieke prestatiedata te koppelen aan wedstrijdresultaten wordt nagegaan welke fysieke kenmerken statistisch samenhangen met sportief succes. 
De doelgroep van dit onderzoek bestaat uit performance analisten en scouts binnen het professionele voetbal die nood hebben aan objectieve, interpreteerbare en data-gedreven ondersteuning bij scoutingbeslissingen.
\newline\newline
De centrale onderzoeksvraag van deze bachelorproef luidt: \textbf{In welke mate kunnen fysieke spelersstatistieken, geanalyseerd met explainable machine learning technieken, bijdragen aan het identificeren van relevante spelersprofielen voor voetbal scouting?}
\newline
Het doel van dit onderzoek is het ontwikkelen van een proof-of-concept dat aantoont hoe machine learning en explainable AI kunnen worden ingezet om fysieke prestatiedata te vertalen naar bruikbare inzichten voor scouting.

Enkele deelvragen met betrekking tot het probleemdomein zijn: 
\begin{itemize}
    \item Welke fysieke spelersattributen worden binnen professioneel voetbal relevant geacht voor prestaties en scouting?
    \item In welke mate vertonen fysieke spelersstatistieken een verband met wedstrijdresultaten?
    \item Welke contextuele factoren (positie, thuis/uit, competitieniveau) beïnvloeden de relevantie van fysieke attributen?
\end{itemize}

Enkele deelvragen met betrekking tot het oplossingsdomein zijn: 
\begin{itemize}
    \item Welke machine learning modellen zijn geschikt voor het voorspellen van wedstrijdresultaten op basis van fysieke data?
    \item Hoe kunnen explainable machine learning technieken worden toegepast om de bijdrage van fysieke attributen te interpreteren?
    \item Hoe kunnen deze inzichten vertaald worden naar praktische toepassingen binnen voetbal scouting?
\end{itemize}

%---------- Stand van zaken ---------------------------------------------------

\section{Stand van zaken}%
\label{sec:literatuurstudie}

De toepassing van data-analyse en machine learning binnen de sportwereld is de afgelopen jaren sterk toegenomen. In het bijzonder binnen professioneel voetbal worden steeds grotere hoeveelheden data verzameld, gaande van wedstrijd- en eventdata tot gedetailleerde fysieke prestatiedata afkomstig van GPS-systemen en wearables. Deze evolutie heeft geleid tot een groeiend onderzoeksdomein waarin machine learning technieken worden ingezet om sportprestaties te analyseren en te voorspellen.

Een recente systematische review van \textcite{Pietraszewski2025_AI_SportsAnalytics} toont aan dat artificial intelligence en machine learning een steeds prominentere rol spelen binnen sport analytics. Volgens de auteurs worden deze technieken voornamelijk ingezet voor prestatie-evaluatie, wedstrijdvoorspelling en blessurepreventie. Hoewel het merendeel van de studies focust op technische en tactische indicatoren, benadrukt de review dat ook fysieke prestatiedata steeds vaker wordt opgenomen in voorspellende modellen. Dit onderstreept het potentieel van machine learning voor het analyseren van fysieke spelersattributen.

Binnen voetbalcontext wordt machine learning voornamelijk toegepast op wedstrijd- en eventdata. \textcite{Pappalardo2018_PlayeRank} introduceren met het PlayeRank-framework een data-gedreven methode om spelersprestaties te evalueren en te rangschikken. Hoewel dit framework voornamelijk technische en tactische statistieken gebruikt, illustreert het onderzoek hoe machine learning kan worden ingezet om multidimensionale spelersdata te vertalen naar objectieve evaluaties. Deze aanpak vormt een belangrijke inspiratie voor het toepassen van data-gedreven methoden binnen scoutingprocessen.

Naast technische statistieken is er toenemende aandacht voor het gebruik van fysieke prestatiedata binnen machine learning modellen. \textcite{Mandorino2023_PlayerReadiness} tonen aan dat fysieke parameters zoals acceleraties, afstanden en intensieve acties succesvol kunnen worden geïntegreerd in machine learning modellen om de fysieke paraatheid van voetballers te kwantificeren. Dit bevestigt dat fysieke spelersstatistieken niet enkel relevant zijn voor load monitoring, maar ook waardevolle input vormen voor analytische modellen.

Machine learning wordt ook frequent ingezet voor het voorspellen van wedstrijdresultaten. In een studie binnen het \textit{Journal of Big Data} tonen \textcite{BigData2024_SoccerOutcomePrediction} aan dat zowel klassieke machine learning modellen als meer geavanceerde technieken in staat zijn om voetbalwedstrijden te classificeren op basis van wedstrijdgerelateerde features. Hoewel fysieke statistieken in deze studie niet centraal staan, bevestigt het onderzoek dat match outcome prediction een geschikt kader vormt om de impact van specifieke featuregroepen, zoals fysieke teamkenmerken, te analyseren.

Een belangrijk aandachtspunt binnen sport analytics is de interpreteerbaarheid van machine learning modellen. In professionele omgevingen, zoals voetbal scouting, is het essentieel dat beslissingen niet enkel gebaseerd zijn op voorspellende prestaties, maar ook begrijpelijk en verantwoordbaar zijn. \textcite{CalderonDiaz2023_ExplainableML_Soccer} tonen aan dat explainable machine learning technieken, zoals feature importance en SHAP-waarden, inzicht kunnen bieden in de bijdrage van individuele fysieke kenmerken aan modelvoorspellingen. Hoewel hun onderzoek gericht is op blessurepredictie, is de gebruikte methodologie rechtstreeks toepasbaar op prestatie- en scoutinganalyses.

Tot slot benadrukken methodologische studies het belang van correcte dataverwerking en feature engineering binnen sport analytics. \textcite{Springer2024_MethodologyChallenges} bespreken verschillende uitdagingen met betrekking tot datakwaliteit, aggregatie en evaluatie van machine learning modellen in sportcontexten. Deze uitdagingen zijn bijzonder relevant wanneer fysieke spelersdata wordt geaggregeerd naar teamniveau, zoals vereist is binnen wedstrijdanalyse en scoutinggerichte toepassingen.

Samenvattend toont de literatuur aan dat machine learning een waardevol instrument is binnen sport analytics, maar dat de toepassing van fysieke prestatiedata binnen scouting nog beperkt onderzocht is. Daarnaast blijkt interpreteerbaarheid een cruciale factor voor praktische toepasbaarheid. Deze bachelorproef sluit aan bij deze vaststellingen door fysieke spelersattributen centraal te stellen en explainable machine learning technieken toe te passen binnen een scoutingcontext.


% Voor literatuurverwijzingen zijn er twee belangrijke commando's:
% \autocite{KEY} => (Auteur, jaartal) Gebruik dit als de naam van de auteur
%   geen onderdeel is van de zin.
% \textcite{KEY} => Auteur (jaartal)  Gebruik dit als de auteursnaam wel een
%   functie heeft in de zin (bv. ``Uit onderzoek door Doll & Hill (1954) bleek
%   ...'')


%---------- Methodologie ------------------------------------------------------
\section{Methodologie}%
\label{sec:methodologie}

Het onderzoek wordt uitgevoerd volgens een gestructureerde data science aanpak en bestaat uit meerdere fasen.

In een eerste fase wordt via literatuurstudie onderzocht welke fysieke spelersattributen relevant zijn binnen professioneel voetbal en hoe machine learning reeds wordt toegepast binnen sport analytics. Op basis hiervan wordt een selectie gemaakt van fysieke prestatie-indicatoren.

In de tweede fase wordt een dataset samengesteld bestaande uit wedstrijdgegevens, fysieke spelersstatistieken en wedstrijdresultaten over meerdere seizoenen. Individuele spelersstatistieken worden via feature engineering geaggregeerd naar teamniveau, rekening houdend met positie en speelminuten.

Vervolgens worden verschillende classificatiemodellen ontwikkeld om wedstrijdresultaten (winst, gelijkspel, verlies) te voorspellen op basis van deze fysieke teamfeatures. Hierbij wordt gestart met een interpreteerbaar basismodel, gevolgd door complexere machine learning modellen.

In de laatste fase worden explainable machine learning technieken toegepast om de bijdrage van fysieke attributen aan de modelvoorspellingen te analyseren. De verkregen inzichten worden geïnterpreteerd in functie van voetbal scouting en vertaald naar concrete aanbevelingen.

%---------- Verwachte resultaten ----------------------------------------------
\section{Verwacht resultaat, conclusie}%
\label{sec:verwachte_resultaten}

Er wordt verwacht dat dit onderzoek aantoont dat bepaalde fysieke spelersattributen een significante bijdrage leveren aan wedstrijdresultaten en dat deze attributen op een objectieve manier kunnen worden geïdentificeerd met behulp van machine learning. Daarnaast wordt verwacht dat explainable AI technieken toelaten om deze resultaten op een transparante en interpreteerbare manier te communiceren.

De bachelorproef zal aantonen dat fysieke prestatiedata niet enkel waardevol is voor trainingsopvolging, maar ook een relevante rol kan spelen binnen voetbal scouting. Het ontwikkelde proof-of-concept biedt een eerste stap richting data-gedreven scoutingbeslissingen waarbij objectieve fysieke profielen worden gecombineerd met menselijke expertise.